\appendices
\section{Comparaciones detalladas suplementarias}
Las tablas conservan escenarios secundarios sin sustituir el criterio principal. La Tabla~\ref{tab:stress} contiene todas las familias bajo mecanismos añadidos; la Tabla~\ref{tab:variation} distingue variación de medias por semilla y por máscara. La desviación entre semillas no se define para una referencia determinista con un solo ajuste; no se sustituye silenciosamente por cero. Son resúmenes condicionados, no componentes independientes de varianza poblacional.
\begin{table*}[!t]
\caption{MASE descriptivo bajo pérdidas correlacionadas, prolongadas y MAR dependiente del retardo. La etiqueta nominal no equivale siempre a tasa realizada.}\label{tab:stress}
\centering\footnotesize
\FitTable{\begin{tabular}{lrrrrrr}
\toprule
Modelo & CPU corr. & CPU larga & CPU lag-MAR & Mem. corr. & Mem. larga & Mem. lag-MAR \\
\midrule
Persistencia & 1.0836 & 1.2113 & 1.0934 & 1.2422 & 1.4198 & 1.2493 \\
Media móvil & 1.0326 & 1.1509 & 1.0521 & 1.1885 & 1.3576 & 1.2091 \\
Suavizamiento exponencial & 1.0332 & 1.1470 & 1.0522 & 1.2024 & 1.3700 & 1.2254 \\
Ridge & 1.2453 & 1.3294 & 1.2518 & 1.3217 & 1.4826 & 1.3355 \\
DLinear & 1.1235 & 1.2222 & 1.1326 & 1.1956 & 1.3701 & 1.2120 \\
NLinear & 1.0121 & 1.1375 & 1.0281 & 1.1737 & 1.3532 & 1.1902 \\
E5 & 1.1194 & 1.2187 & 1.1300 & 1.2250 & 1.4072 & 1.2452 \\
A0 & 1.1300 & 1.2285 & 1.1393 & 1.1881 & 1.3638 & 1.2048 \\
A1 & 1.1285 & 1.2273 & 1.1377 & 1.1942 & 1.3701 & 1.2109 \\
A2 & 1.1299 & 1.2289 & 1.1393 & 1.1985 & 1.3783 & 1.2151 \\
A3 & 1.1342 & 1.2330 & 1.1428 & 1.1932 & 1.3743 & 1.2106 \\
GRU (propia) & 1.1879 & 1.3095 & 1.2019 & 1.2932 & 1.5458 & 1.3110 \\
GRU (grafo) & 1.1603 & 1.2816 & 1.1761 & 1.4235 & 1.7076 & 1.4393 \\
GRU-D adaptada & 1.1394 & 1.2605 & 1.1592 & 1.3076 & 1.5114 & 1.3286 \\
PatchTST (pequeño) & 1.0053 & 1.1323 & 1.0223 & 1.1712 & 1.3485 & 1.1894 \\
\bottomrule
\end{tabular}
}
\end{table*}
\begin{table*}[!t]
\caption{Desviaciones de medias por semilla y máscara, principal con pérdidas. No son incertidumbres de infraestructuras independientes.}\label{tab:variation}
\centering\footnotesize
\FitTable{\begin{tabular}{lrrrr}
\toprule
Modelo & CPU DE semilla & CPU DE máscara & Memoria DE semilla & Memoria DE máscara \\
\midrule
Persistencia & sin definir & 0.0033 & sin definir & 0.0065 \\
Media móvil & sin definir & 0.0031 & sin definir & 0.0062 \\
Suavizamiento exponencial & sin definir & 0.0034 & sin definir & 0.0082 \\
Ridge & sin definir & 0.0021 & sin definir & 0.0063 \\
DLinear & 0.0332 & 0.0023 & 0.0161 & 0.0062 \\
NLinear & 0.0081 & 0.0030 & 0.0112 & 0.0062 \\
E5 & 0.0227 & 0.0025 & 0.0179 & 0.0066 \\
A0 & 0.0265 & 0.0024 & 0.0126 & 0.0063 \\
A1 & 0.0238 & 0.0024 & 0.0215 & 0.0063 \\
A2 & 0.0262 & 0.0024 & 0.0179 & 0.0063 \\
A3 & 0.0222 & 0.0024 & 0.0154 & 0.0063 \\
GRU (propia) & 0.0634 & 0.0018 & 0.0341 & 0.0050 \\
GRU (grafo) & 0.0313 & 0.0022 & 0.0403 & 0.0051 \\
GRU-D adaptada & 0.0426 & 0.0020 & 0.0463 & 0.0050 \\
PatchTST (pequeño) & 0.0093 & 0.0030 & 0.0105 & 0.0062 \\
\bottomrule
\end{tabular}
}
\end{table*}
\begin{table}[!t]
\caption{Intervalos95\% cruzados por semilla, máscara y bloques circulares para referencia menos A3. Los bloques24/48/96 equivalen a12/24/48 horas nominales. No se afirma cobertura predictiva o poblacional.}\label{tab:block}
\centering\footnotesize
\FitTable{\begin{tabular}{lrrrr}
\toprule
Recurso & Referencia & Pasos del bloque & Inferior & Superior \\
\midrule
CPU & DLinear & 24 & -0.02818 & 0.00581 \\
CPU & DLinear & 48 & -0.02732 & 0.00470 \\
CPU & DLinear & 96 & -0.02944 & 0.00518 \\
CPU & E5 & 24 & -0.02903 & -0.00094 \\
CPU & E5 & 48 & -0.02977 & -0.00120 \\
CPU & E5 & 96 & -0.02966 & -0.00011 \\
Memoria & DLinear & 24 & -0.01024 & 0.01019 \\
Memoria & DLinear & 48 & -0.00959 & 0.01035 \\
Memoria & DLinear & 96 & -0.00988 & 0.00981 \\
Memoria & E5 & 24 & 0.01264 & 0.03961 \\
Memoria & E5 & 48 & 0.01312 & 0.03981 \\
Memoria & E5 & 96 & 0.01233 & 0.03971 \\
\bottomrule
\end{tabular}
}
\end{table}
\begin{table*}[!t]
\caption{MASE con objetivos comunes: CPU a la izquierda y memoria a la derecha, horizontes en minutos. Tres semillas, cinco máscaras y seis condiciones con pérdidas. H30 es un subconjunto del principal.}\label{tab:horizon}
\centering\footnotesize

\setlength{\tabcolsep}{4pt}
\begin{minipage}[t]{.485\textwidth}
\centering\FitTable{\begin{tabular}{lrrrrr}
\toprule
Modelo & 30 & 60 & 120 & 180 & 360 \\
\midrule
Persistencia & 1.1159 & 1.1381 & 1.3408 & 1.4851 & 1.7819 \\
Media móvil & 1.0666 & 1.1420 & 1.3589 & 1.4831 & 1.7010 \\
Suavizamiento exponencial & 1.0686 & 1.1469 & 1.3009 & 1.4242 & 1.6651 \\
Ridge & 1.2672 & 1.3856 & 1.6037 & 1.7441 & 2.0420 \\
DLinear & 1.1453 & 1.2419 & 1.4287 & 1.6340 & 1.9346 \\
E5 & 1.1384 & 1.2541 & 1.4506 & 1.6121 & 1.9251 \\
A3 & 1.1584 & 1.2492 & 1.4448 & 1.6451 & 1.9779 \\
PatchTST (pequeño) & 1.0445 & 1.1217 & 1.2814 & 1.4077 & 1.6399 \\
\bottomrule
\end{tabular}
}
\end{minipage}
\hfill
\begin{minipage}[t]{.485\textwidth}
\centering\FitTable{\begin{tabular}{lrrrrr}
\toprule
Modelo & 30 & 60 & 120 & 180 & 360 \\
\midrule
Persistencia & 1.2862 & 1.3398 & 1.6072 & 1.8097 & 2.2313 \\
Media móvil & 1.2364 & 1.3384 & 1.6286 & 1.7971 & 2.2460 \\
Suavizamiento exponencial & 1.2544 & 1.3549 & 1.5611 & 1.7315 & 2.1303 \\
Ridge & 1.3640 & 1.5179 & 1.8084 & 2.0411 & 2.5289 \\
DLinear & 1.2528 & 1.3958 & 1.6037 & 1.8014 & 2.3050 \\
E5 & 1.2680 & 1.4250 & 1.6459 & 1.8648 & 2.3222 \\
A3 & 1.2436 & 1.3474 & 1.6060 & 1.8353 & 2.3141 \\
PatchTST (pequeño) & 1.2264 & 1.3377 & 1.5866 & 1.7572 & 2.1681 \\
\bottomrule
\end{tabular}
}
\end{minipage}
\end{table*}
\begin{table*}[!t]
\caption{Historias semihorarias con objetivos comunes y horizonte treinta minutos (CPU a la izquierda y memoria a la derecha). Diseño estrella descriptivo.}\label{tab:history}
\centering\footnotesize

\setlength{\tabcolsep}{4pt}
\begin{minipage}[t]{.485\textwidth}
\centering\FitTable{\begin{tabular}{lrrr}
\toprule
Modelo & 12 & 24 & 48 \\
\midrule
Persistencia & 1.1159 & 1.1159 & 1.1159 \\
Media móvil & 1.0666 & 1.0666 & 1.0666 \\
Suavizamiento exponencial & 1.0686 & 1.0625 & 1.0618 \\
Ridge & 1.2672 & 1.2403 & 1.1504 \\
DLinear & 1.1453 & 1.1372 & 1.1072 \\
E5 & 1.1384 & 1.1402 & 1.1071 \\
A3 & 1.1584 & 1.1572 & 1.0966 \\
PatchTST (pequeño) & 1.0445 & 1.0421 & 1.0618 \\
\bottomrule
\end{tabular}
}
\end{minipage}
\hfill
\begin{minipage}[t]{.485\textwidth}
\centering\FitTable{\begin{tabular}{lrrr}
\toprule
Modelo & 12 & 24 & 48 \\
\midrule
Persistencia & 1.2862 & 1.2862 & 1.2862 \\
Media móvil & 1.2364 & 1.2364 & 1.2364 \\
Suavizamiento exponencial & 1.2544 & 1.2501 & 1.2494 \\
Ridge & 1.3640 & 1.3385 & 1.3063 \\
DLinear & 1.2528 & 1.2543 & 1.2610 \\
E5 & 1.2680 & 1.2680 & 1.2681 \\
A3 & 1.2436 & 1.2521 & 1.2560 \\
PatchTST (pequeño) & 1.2264 & 1.2190 & 1.2413 \\
\bottomrule
\end{tabular}
}
\end{minipage}
\end{table*}
\begin{table*}[!t]
\caption{Transferencia Alibaba, seis condiciones con pérdidas. MAE/RMSE absolutos utilizan normalización publicada de ambos recursos, no núcleos o bytes.}\label{tab:external}
\centering\footnotesize
\FitTable{\begin{tabular}{lrrrrrr}
\toprule
Modelo & CPU MASE & CPU MAE & CPU RMSE & Mem. MASE & Mem. MAE & Mem. RMSE \\
\midrule
Persistencia & 2.3795 & 0.007548 & 0.011324 & 7.5150 & 0.001085 & 0.002302 \\
Media móvil & 2.8906 & 0.009822 & 0.013787 & 9.6034 & 0.001392 & 0.002726 \\
Suavizamiento exponencial & 3.3095 & 0.011482 & 0.015495 & 10.2802 & 0.001847 & 0.003271 \\
Ridge & 3.6334 & 0.012690 & 0.016598 & 10.0263 & 0.001885 & 0.003208 \\
DLinear & 3.3262 & 0.011559 & 0.015422 & 9.8853 & 0.001762 & 0.003090 \\
NLinear & 3.0566 & 0.010455 & 0.014191 & 9.5989 & 0.001658 & 0.002972 \\
E5 & 3.3842 & 0.011791 & 0.015759 & 9.7985 & 0.001767 & 0.003104 \\
A0 & 3.3521 & 0.011664 & 0.015584 & 9.9413 & 0.001760 & 0.003097 \\
A1 & 3.3481 & 0.011650 & 0.015572 & 9.9291 & 0.001759 & 0.003095 \\
A2 & 3.3520 & 0.011663 & 0.015584 & 9.9668 & 0.001780 & 0.003117 \\
A3 & 3.3593 & 0.011686 & 0.015601 & 9.9535 & 0.001768 & 0.003105 \\
GRU (propia) & 5.8743 & 0.020136 & 0.024363 & 17.6403 & 0.004010 & 0.006199 \\
GRU-D adaptada & 5.2232 & 0.017836 & 0.022202 & 17.5279 & 0.003966 & 0.006144 \\
PatchTST (pequeño) & 3.1112 & 0.010653 & 0.015257 & 10.5876 & 0.001559 & 0.002796 \\
\bottomrule
\end{tabular}
}
\end{table*}
La Tabla~\ref{tab:historical} conserva los promedios históricos con telemetría completa incluida. La Tabla~\ref{tab:matrix} documenta el alcance de la ampliación y la repetición de filas de métricas.
\begin{table}[!t]
\caption{Resumen histórico de las siete condiciones originales, incluida telemetría completa. Su alcance difiere de la tabla histórica principal de seis condiciones con pérdidas. V2 es alias y el extremo GinAR se conserva.}\label{tab:historical}
\centering\footnotesize
\FitTable{\begin{tabular}{lrr}
\toprule
Modeloo histórico & CPU MASE & Memoria MASE \\
\midrule
persistencia & 1.0761 & 1.3112 \\
E0 & 0.9874 & 1.1283 \\
GRU\_propia & 0.9495 & 1.1580 \\
GRU\_llamadas & 0.9324 & 1.1893 \\
GinAR & 704115.0298 & 2.5812 \\
BiTGraph\_codigo\_autor & 0.9914 & 1.2009 \\
E1 & 1.0179 & 1.1283 \\
E2 & 1.0170 & 1.1178 \\
E3 & 0.9994 & 1.1267 \\
E4 & 0.9852 & 1.1283 \\
E5 & 0.9979 & 1.1136 \\
C0 & 1.0170 & 1.1283 \\
C1 & 1.0168 & 1.1278 \\
C2 & 1.0171 & 1.1284 \\
C3 & 1.0168 & 1.1280 \\
V2 & 1.0168 & 1.1278 \\
\bottomrule
\end{tabular}
}
\end{table}
\begin{table}[!t]
\caption{Matriz de métricas ejecutadas. Las filas repiten resúmenes por modelo, semilla, máscara, condición y servicio; no son datasets independientes.}\label{tab:matrix}
\centering\footnotesize
\FitTable{\begin{tabular}{lrrrrr}
\toprule
Conjunto & L & h pasos & Modeloos & Máscaras & Filas de métricas \\
\midrule
alibaba\_transfer & 12 & 1 & 14 & 20 & 1331200 \\
new\_services & 12 & 1 & 15 & 20 & 1459200 \\
new\_services & 12 & 2 & 8 & 5 & 51200 \\
new\_services & 12 & 4 & 8 & 5 & 51200 \\
new\_services & 12 & 6 & 8 & 5 & 51200 \\
new\_services & 12 & 12 & 8 & 5 & 51200 \\
new\_services & 24 & 1 & 8 & 5 & 51200 \\
new\_services & 48 & 1 & 8 & 5 & 51200 \\
\bottomrule
\end{tabular}
}
\end{table}
\section{Reconstrucción y procedencia}
El manifiesto de predicciones registra cada archivo nuevo y sus hashes de archivo y arreglos; el de checkpoints registra pesos y huellas del predictor. El archivo de métricas conserva precisión completa y semántica del objetivo. El redondeo de tablas no modifica las fuentes. La trazabilidad del artículo relaciona cifras, filtros y representaciones con fuentes verificadas. Los hashes de código del protocolo se comprobaron nuevamente durante análisis.

El analizador reconstruye en float64 las ecuaciones, cantidades de objetivos e índices desde predicciones guardadas. Regenerar máscaras con el generador congelado es una comprobación de consistencia, no un segundo algoritmo independiente. Las pruebas cubren límites/huecos, invariancia de prefijos futuros MAR, perturbación oculta con máscara fija, procedencia, modelos finitos CPU/CUDA e inicialización afín. Los ensayos del mismo entorno verifican recuperación y reproducción CUDA. No demuestran ausencia universal de fugas en datos del propietario.

La imagen Docker derivada reutiliza una base local. Sus identificadores y dependencias están registrados, sin certificar reconstrucción limpia en otra máquina. Se corrigió un defecto de inicialización del muestreo CPU del monitor: RSS previo sigue válido y ceros CPU previos se excluyen. Se preserva la enmienda. Épocas, RSS, asignación PyTorch y telemetría global mantienen definiciones separadas. El README describe reconstrucción limitada sin volver a seleccionar ni entrenar.

\subsection{Mediciones separadas de memoria y tiempos históricos}
La Tabla~\ref{tab:historical-bench} conserva mediciones seleccionadas CUDA B128 sobre fuente28. La Tabla~\ref{tab:memory-bench} muestra inferencia CUDA B1 en destino32 para el objetivo CPU. RSS describe memoria del proceso y el pico CUDA corresponde al asignador PyTorch. El contexto y las memorias de partida difieren entre protocolos. No se convierten directamente en tamaño de parámetros.
\begin{table}[!t]
\caption{Histórico B128 CUDA directo, semilla101. RSS es pico del proceso reportado y CUDA pico de memoria asignada. Los segundos incluyen ajuste y validación.}\label{tab:historical-bench}
\centering\footnotesize
\FitTable{\begin{tabular}{lrrrrr}
\toprule
Modelo / recurso & p50 (ms) & p95 (ms) & RSS (MiB) & CUDA (MiB) & s \\
\midrule
E0 / Memoria & 0.287 & 0.602 & 840.2 & 64.84 & 7.805 \\
E2 / Memoria & 0.595 & 0.692 & 849.6 & 64.21 & 5.199 \\
E4 / Memoria & 0.441 & 0.937 & 889.8 & 64.51 & 5.295 \\
E5 / Memoria & 0.801 & 0.869 & 930.7 & 77.35 & 6.342 \\
V2 / Memoria & 0.416 & 0.556 & 914.1 & 64.51 & 4.630 \\
E0 / CPU & 0.284 & 0.483 & 839.7 & 64.84 & 12.799 \\
E2 / CPU & 0.585 & 0.681 & 848.7 & 64.21 & 9.521 \\
E4 / CPU & 0.432 & 0.675 & 889.2 & 64.51 & 15.421 \\
E5 / CPU & 0.814 & 0.936 & 930.8 & 77.35 & 7.473 \\
V2 / CPU & 0.691 & 0.767 & 860.4 & 64.21 & 5.999 \\
\bottomrule
\end{tabular}
}
\end{table}
\begin{table}[!t]
\caption{CUDA B1 ampliado con32 servicios. RSS observado y pico del asignador CUDA en MiB. GRU-D utiliza reproducción verificada. Los totales del dispositivo se informan por separado.}\label{tab:memory-bench}
\centering\footnotesize
\FitTable{\begin{tabular}{lrr}
\toprule
Modelo & RSS (MiB) & CUDA (MiB) \\
\midrule
A0 & 1160.49 & 114.861 \\
A1 & 1157.96 & 114.862 \\
A2 & 1160.36 & 114.861 \\
A3 & 1160.38 & 114.862 \\
DLinear & 1160.27 & 114.866 \\
E5 & 1157.80 & 114.877 \\
GRU-D adaptada & 1136.98 & 114.875 \\
GRU (grafo) & 1160.43 & 123.216 \\
GRU (propia) & 1157.50 & 123.216 \\
Media móvil & 1137.01 & 114.865 \\
NLinear & 1157.62 & 114.861 \\
PatchTST (pequeño) & 1096.40 & 82.992 \\
Persistencia & 1137.00 & 114.859 \\
Ridge & 1157.65 & 114.861 \\
Suavizamiento exponencial & 1160.39 & 114.860 \\
\bottomrule
\end{tabular}
}
\end{table}
